% ═════════════════════════════════════════════════════════════
% A2 --- QR Factorization
\label{sec:qr}
% ═════════════════════════════════════════════════════════════

\section{QR Solution of the Regularized Least-Squares Problem}
\label{sec:qr_solution}

\noindent We now turn from the optimization point of view to the direct linear algebra point 
of view. The same regularized least squares problem can be solved without generating 
an iterative sequence of approximations: one factorizes the augmented matrix

\begin{equation}
    A \coloneqq \Xhat =
    \begin{bmatrix}
        X^T \\ \lambda I_m
    \end{bmatrix}
    \in \R^{(n+m)\times m},
    \qquad
    b \coloneqq \yhat =
    \begin{bmatrix}
        y \\ 0
    \end{bmatrix}
    \in \R^{n+m},
\end{equation}

\noindent and then solves the resulting triangular problem. Since $\lambda > 0$, 
Lemma~\ref{lemma:fullrank} shows that $A$ has full column rank. Therefore the 
least squares problem

\begin{equation}
    \min_{w\in\R^m} \norm{Aw-b}
\end{equation}

\noindent has a unique minimizer.

\subsection{Thin QR Factorization}

\noindent For a full column rank matrix $A\in\R^{p\times m}$, with $p=n+m\geq m$, the thin QR 
factorization is

\begin{equation}
\label{eq:thin_qr}
    A = Q_1 R,
    \qquad
    Q_1 \in \R^{p\times m}, \quad Q_1^T Q_1 = I_m,
    \quad
    R \in \R^{m\times m},
\end{equation}

\noindent where $R$ is upper triangular and nonsingular.


\noindent Let $Q_2 \in \R^{p\times n}$ be a matrix whose columns form an orthonormal basis of the orthogonal
complement of $\operatorname{col}(Q_1)$. Then

\begin{equation}
    Q =
    \begin{bmatrix}
    Q_1 & Q_2
    \end{bmatrix}
    \in \R^{p\times p}
\end{equation}

\noindent is an orthogonal matrix:

\begin{equation}
    Q^T Q = Q Q^T = I_p.
\end{equation}

\noindent Therefore,

\begin{equation}
    Q_1^T Q_1 = I_m,
    \qquad
    Q_2^T Q_2 = I_{n},
    \qquad
    Q_2^T Q_1 = 0.
\end{equation}

\noindent Hence,

\begin{equation}
    Q^T A = Q^T Q_1 R =
    \begin{bmatrix}
        R \\ 0
    \end{bmatrix},
    \qquad
    Q^T b =
    \begin{bmatrix}
        Q_1^T b \\ Q_2^T b
    \end{bmatrix} =
    \begin{bmatrix}
        c \\ d
    \end{bmatrix},
    \quad c\in\R^m,\ d\in\R^{n}.
\end{equation}

\noindent Orthogonal transformations preserve the Euclidean norm, hence

\begin{equation}
\label{eq:qr_residual_split}
    \norm{Aw-b}^2
    =
    \norm{Q^T(Aw-b)}^2
    =
    \norm{Rw-c}^2 + \norm{d}^2.
\end{equation}

\noindent Since the term $\norm{d}^2$ is independent of $w$, the minimization problem reduces to
\begin{equation}
\label{eq:qr_triangular_solve}
    \arg\min_{w\in\R^m} \norm{Aw-b}^2
    =
    \arg\min_{w\in\R^m} \norm{Rw-c}^2.
\end{equation}
\noindent Since $R$ is nonsingular, there exists a unique $w$ such that
\begin{equation}
    Rw-c=0,
    \qquad\text{that is,}\qquad
    Rw=c=Q_1^T b.
\end{equation}

\noindent Since $A^T A = (Q_1R)^{T}(Q_1R)$, this is the QR analogue of solving the normal equations. Indeed,

\begin{equation}
\label{eq:qr_normal_equivalence}
    R^T R = A^T A = XX^T + \lambda^2 I_m,
    \qquad
    R^T c = A^T b = Xy.
\end{equation}

\noindent Thus QR and the normal equations define the same solution in exact arithmetic. Numerically, however, 
QR is preferable because it avoids explicitly forming and solving with $A^T A$, whose 
condition number is squared:

\begin{equation}
\label{eq:qr_condition_squaring}
    \kappa_2(A^T A) = \kappa_2(A)^{2}.
\end{equation}

\noindent If $X\neq0$, this distinction is important as $\lambda\to0^+$, since
Eq.~\eqref{eq:condition} gives
$\kappa_2(A)\sim\sigma_{\max}(X)/\lambda$.

\subsection{Algorithmic Structure}

\noindent The QR-based solver can be summarized as follows:

\begin{algorithm}[H]
\caption{QR Solver for the Regularized Least Squares Problem}
\label{alg:qr_solver}
\begin{algorithmic}[1]
    \Require Data matrix $X\in\R^{m\times n}$; response $y\in\R^n$; parameter $\lambda>0$
    \State Form the augmented matrix $A = [X^T;\lambda I_m]$ and vector $b=[y;0]$
    \State Compute a thin QR factorization $A=Q_1R$
    \State Compute $c \gets Q_1^T b$
    \State Solve $Rw=c$ by back-substitution
    \State \Return $w$
\end{algorithmic}
\end{algorithm}

\noindent In practice, the matrix $Q_1$ is not formed explicitly. Forming it would require 
$O((n+m)m)$ storage and applying it as a dense matrix would obscure the sparse and 
structured nature of $A$. Instead, the orthogonal transformations used during the 
factorization are stored in compact form and later applied to $b$.

\subsection{Cost for the QR Solver}
\label{sec:qr_cost_project}
\noindent For a generic dense matrix $A\in\R^{p\times m}$, with $p=n+m\geq m$, computing a thin QR factorization requires

\begin{equation}
    O(pm^2) = O((n+m)m^2)
\end{equation}

\noindent Once the factorization is available, computing the transformed right-hand side $c=Q_1^T b$ requires at most $O(pm)$ additional operations, whereas solving the upper triangular system

\begin{equation}
    Rw=c
\end{equation}

\noindent by back-substitution costs $O(m^2)$.
Therefore, the factorization dominates the overall computational cost. Since $p=n+m$, treating the augmented matrix $A$ as a generic dense matrix gives

\begin{equation}
    O((n+m)m^2) = O(nm^2+m^3),
\end{equation}

\noindent This bound contains a cubic term in $m$.
This is the baseline complexity of the QR solver before exploiting the particular block structure of the augmented matrix. The structured implementation developed in Section~\ref{sec:householder_qr} reduces this cost substantially.


% ═════════════════════════════════════════════════════════════
\section{QR with Householder Reflectors}
\label{sec:householder_qr}
% ═════════════════════════════════════════════════════════════

\noindent Householder reflectors are the standard stable mechanism for computing QR factorizations~\cite{trefethen1997,golub2013}. The implementation stores a normalized reflector vector $u$, so a Householder reflector has the form

\begin{equation}
\label{eq:householder_def}
    H = I - 2uu^T,
    \qquad
    \norm{u}_2=1.
\end{equation}

\noindent Equivalently, for an unnormalized nonzero vector $v$, one may write
$H=I-\tau vv^T$ with $\tau=2/(v^T v)$. In either representation, $H$ is
symmetric and orthogonal:

\begin{equation}
    H^T = H,
    \qquad
    H^T H = I.
\end{equation}

\noindent Given a nonzero vector $z\in\R^\ell$, the reflector is chosen so that all
components below the first one are annihilated. To match the implementation and
avoid overflow or underflow in the norm calculation, define

\begin{equation}
\label{eq:stable_householder_vector}
    s=\norm{z}_\infty,
    \qquad
    \bar z=\frac{z}{s},
    \qquad
    \nu=\norm{\bar z}_2,
    \qquad
    \sigma=\begin{cases}1,&z_1\geq0,\\-1,&z_1<0,
    \end{cases}
\end{equation}

\noindent and then

\begin{equation}
    \alpha=-\sigma s\nu,
    \qquad
    v=\bar z+\sigma\nu e_1,
    \qquad
    u=\frac{v}{\norm{v}_2}.
\end{equation}

\noindent With $H$ as in Eq.~\eqref{eq:householder_def}, this construction gives

\begin{equation}
    Hz = \alpha e_1
    = 
    \begin{bmatrix}
        -\operatorname{sign}(z_1)\norm{z}_2 \\
        0 \\
        \vdots \\
        0
    \end{bmatrix}.
\end{equation}

\noindent Thus the convention is $\operatorname{sign}(0)=1$. The sign choice makes
the first component of $v$ an addition of like-signed quantities, thereby
avoiding subtractive cancellation. If $z=0$, the code returns $u=0$ and
$\alpha=0$ as an identity-reflector sentinel; reflector-application routines
skip this sentinel.

\noindent At step $i$, the current active column is therefore mapped to a vector
whose entries below position $i$ are zero, while its $i$-th entry becomes
$\alpha$.

\subsection{Dense Householder QR}

\noindent The function \texttt{naive\_qr\_solver} applies Householder QR to the
explicit augmented matrix as if it were generic and dense; here \emph{naive}
means that the block structure is ignored, not that the transformation is
unstable. For $A\in\R^{p\times m}$, the factorization proceeds column by
column. At step $k$, one constructs an active reflector
$\widehat H_k\in\R^{(p-k+1)\times(p-k+1)}$ such that

\begin{equation}
    \widehat H_k A_{k:p,k}
    =
    \begin{bmatrix}
        \alpha_k \\ 0 \\ \vdots \\ 0
    \end{bmatrix}.
\end{equation}

\noindent The same active reflector is applied to the whole trailing submatrix
$A_{k:p,k:m}$. Let $H_k=\operatorname{diag}(I_{k-1},\widehat H_k)$ denote
its extension to $\R^p$. After $m$ steps,

\begin{equation}
    H_m\cdots H_2H_1 A =
    \begin{bmatrix}
        R \\ 0
    \end{bmatrix}.
\end{equation}

\noindent Since each Householder reflector is orthogonal, their product is also
orthogonal. Defining

\begin{equation}
    Q^T = H_m \cdots H_2 H_1,
    \qquad
    Q \in \R^{p\times p},
\end{equation}

\noindent we obtain

\begin{equation}
    Q^T A =
    \begin{bmatrix}
        R \\
        0
    \end{bmatrix},
    \quad
    R \in \R^{m\times m},
\end{equation}

\noindent where $R$ is upper triangular. Multiplying by $Q$ gives

\begin{equation}
    A =
    Q
    \begin{bmatrix}
        R \\
        0
    \end{bmatrix}.
\end{equation}

\noindent Let $Q_1 \in \R^{p\times m}$ denote the first $m$ columns of $Q$. Since the remaining rows of the block matrix above are zero, this reduces to the thin QR factorization

\begin{equation}
    A = Q_1 R,
    \qquad
    Q_1^T Q_1 = I_m.
\end{equation}

\noindent The important implementation point is that $Q$ is never formed explicitly.
The dense implementation first stores the normalized reflector vectors and then
applies them to a copy of $b$, in factorization order:

\begin{equation}
    b_{k:p} \leftarrow \widehat H_k b_{k:p}
    =
    b_{k:p} - 2u_k \left(u_k^T b_{k:p}\right).
\end{equation}

\noindent After all $m$ reflectors have been applied, the work vector contains $Q^T b$. Its first
$m$ entries are

\begin{equation}
    c = Q_1^T b,
\end{equation}

\noindent This is the right-hand side of the triangular system $Rw=c$.

\begin{algorithm}[H]
\caption{Compact Dense Householder QR}
\label{alg:householder_dense}
\begin{algorithmic}[1]
    \Require Matrix $A\in\R^{p\times m}$ with $p\geq m$
    \For{$k=1,\ldots,m$}
        \State $z \gets A_{k:p,k}$
        \State $(u,\alpha) \gets \Call{StableHouseholderVector}{z}$ using Eq.~\eqref{eq:stable_householder_vector}
        \If{$u\neq0$}
            \State $A_{k:p,k:m} \gets A_{k:p,k:m} - 2u(u^T A_{k:p,k:m})$
        \EndIf
        \State $A_{kk}\gets\alpha$; \quad $A_{k+1:p,k}\gets0$
        \State Store $u$
    \EndFor
    \State $R \gets \mathrm{triu}(A_{1:m,1:m})$
    \State \Return $R$ and stored reflectors
\end{algorithmic}
\end{algorithm}

\noindent Algorithm~\ref{alg:householder_dense} is the mathematically canonical version. It is 
useful for explaining QR, but for the project matrix it should be specialized as 
described next.

\subsection{Structure-Based Householder QR}
\label{sec:structured_householder}

\noindent The augmented matrix has the special block structure

\begin{equation}
    A = \begin{bmatrix} X^T \\
    \lambda I_m \end{bmatrix}, \qquad b = \begin{bmatrix} y \\
    0 \end{bmatrix}.
\end{equation}

\noindent The order of the rows is irrelevant for the least-squares problem. Indeed, if $P$ is a permutation matrix, then $P$ is orthogonal and therefore

\begin{equation}
    \norm{Aw-b} = \norm{P(Aw-b)}.
\end{equation}

\noindent We can consequently work with the equivalent permuted system

\begin{equation}
    \label{eq:permuted_augmented} \tilde{A} = \begin{bmatrix} \lambda I_m \\ X^T \end{bmatrix}, \qquad \tilde{b} = \begin{bmatrix} 0 \\ y \end{bmatrix}.
\end{equation}

\noindent This ordering is particularly convenient because the first $m$ rows already form an upper triangular matrix,

\begin{equation}
    R_0=\lambda I_m.
\end{equation}

\noindent The function \texttt{qr\_solver\_structure\_based} performs a column sweep
without forming the full $(m+n)\times m$ matrix. Let $Z=X^T$ denote the dense
lower block. At step $k$, it forms only the active block

\begin{equation}
\label{eq:structure_active_block}
    B_k=
    \begin{bmatrix}
        R_{k,k:m}\\
        Z_{:,k:m}
    \end{bmatrix}
    \in\R^{(n+1)\times(m-k+1)}.
\end{equation}

\noindent A normalized reflector $u_k\in\R^{n+1}$ is computed from the first
column of $B_k$. The update

\begin{equation}
    B_k\leftarrow B_k-2u_k(u_k^T B_k)
\end{equation}

\noindent annihilates all $n$ entries of $Z_{:,k}$ at once. Its first row is
written back to $R_{k,k:m}$ and its remaining rows to $Z_{:,k:m}$. Thus the
algorithm stores $m$ reflectors of fixed length $n+1$, rather than the
decreasing reflectors of length $m+n-k+1$ used by dense QR.

\noindent The same reflectors are applied to the permuted right-hand side without
forming $Q$. At step $k$, only the leading coordinate $c_k$ and the current
dense tail $z\in\R^n$ are active:

\begin{equation}
\label{eq:structure_rhs_update}
    q=\begin{bmatrix}c_k\\z\end{bmatrix},
    \qquad
    q\leftarrow q-2u_k(u_k^T q),
    \qquad
    c_k\leftarrow q_1,
    \quad z\leftarrow q_{2:n+1}.
\end{equation}

\noindent After the last update, $c\in\R^m$ is the leading part of
$Q^T\tilde b$ required by the triangular solve. In the current Python work
buffer, these first $m$ components are returned correctly, but the final
transformed tail $z$ is not copied back; consequently, the last $n$ entries
of the returned buffer must not be interpreted as the complete lower part of
$Q^T\tilde b$.

\begin{algorithm}[H]
\caption{Structure-Based Householder QR by Columns}
\label{alg:structure_based_householder}
\begin{algorithmic}[1]
    \Require $X\in\R^{m\times n}$, $y\in\R^n$, $\lambda>0$
    \State $R\gets\lambda I_m$; \quad $Z\gets X^T$
    \For{$k=1,\ldots,m$}
        \State $B\gets[R_{k,k:m};\,Z_{:,k:m}]$
        \State $(u_k,\alpha_k)\gets\Call{StableHouseholderVector}{B_{:,1}}$
        \If{$u_k\neq0$}
            \State $B\gets B-2u_k(u_k^T B)$
        \EndIf
        \State $B_{11}\gets\alpha_k$; \quad $B_{2:n+1,1}\gets0$
        \State $R_{k,k:m}\gets B_{1,:}$; \quad $Z_{:,k:m}\gets B_{2:n+1,:}$
        \State Store $u_k$
    \EndFor
    \State $c\gets0\in\R^m$; \quad $z\gets y$
    \For{$k=1,\ldots,m$}
        \State $q\gets[c_k;\,z]$
        \If{$u_k\neq0$}
            \State $q\gets q-2u_k(u_k^T q)$
        \EndIf
        \State $c_k\gets q_1$; \quad $z\gets q_{2:n+1}$
    \EndFor
    \State $w\gets\Call{BackSubstitution}{R,c}$
    \State \Return $R,w,c$ and the stored reflectors
\end{algorithmic}
\end{algorithm}

\subsection{Two-Dimensional Householder Row Insertion}
\label{sec:row_insertion_householder}

\noindent A second implementation is based on the mathematically valid idea of
inserting the remaining $n$ rows of $X^T$ one at a time.
At iteration $i$, let $x_i^T = [a_1,\ldots,a_m]$ be the new data row and $y_i$ its corresponding right-hand-side entry. We maintain the current upper triangular factor $R_{i-1}$ and the transformed right-hand side $c_{i-1}$.

\noindent Appending the new row gives

\begin{equation}
    \begin{bmatrix}
        R_{i-1} \\
        x_i^T
    \end{bmatrix}
    =
    \begin{bmatrix}
        R_{11} & R_{12} & R_{13} & \cdots & R_{1m} \\
        0      & R_{22} & R_{23} & \cdots & R_{2m} \\
        0      & 0      & R_{33} & \cdots & R_{3m} \\
        \vdots & \vdots & \ddots & \ddots & \vdots \\
        0      & 0      & \cdots & 0      & R_{mm} \\
        a_1    & a_2    & a_3    & \cdots & a_m
    \end{bmatrix}
    \in \R^{(m+1)\times m}.
\end{equation}

\noindent A sequence of $m$ Householder reflectors is then applied to eliminate the entries of the inserted row one at a time, yielding

\begin{equation}
    \begin{bmatrix}
        R_{i-1} \\
        x_i^T
    \end{bmatrix}
    \longrightarrow
    \begin{bmatrix}
        R_i \\
        0
    \end{bmatrix}
    =
    \begin{bmatrix}
        R^{(i)}_{11} & R^{(i)}_{12} & R^{(i)}_{13} & \cdots & R^{(i)}_{1m} \\
        0            & R^{(i)}_{22} & R^{(i)}_{23} & \cdots & R^{(i)}_{2m} \\
        0            & 0            & R^{(i)}_{33} & \cdots & R^{(i)}_{3m} \\
        \vdots       & \vdots       & \ddots       & \ddots & \vdots \\
        0            & 0            & \cdots       & 0      & R^{(i)}_{mm} \\
        0            & 0            & \cdots       & 0      & 0
    \end{bmatrix}.
\end{equation}

\noindent The same orthogonal transformations must be applied to the corresponding right-hand side:

\begin{equation}
    \begin{bmatrix}
        c_{i-1} \\
        y_i
    \end{bmatrix}
    =
    \begin{bmatrix}
        c_1 \\
        c_2 \\
        \vdots \\
        c_m \\
        y_i
    \end{bmatrix}
    \longrightarrow
    \begin{bmatrix}
        c^{(i)}_1 \\
        c^{(i)}_2 \\
        \vdots \\
        c^{(i)}_m \\
        \eta_i
    \end{bmatrix}
    =
    \begin{bmatrix}
        c_i \\
        \eta_i
    \end{bmatrix}.
\end{equation}

\noindent Unlike the inserted row of the matrix, the last component $\eta_i$ of the transformed right-hand side is not generally zero. Since it corresponds to a zero row in the transformed matrix, it is independent of $w$ and therefore does not affect the least-squares minimizer. Consequently, only the updated vector $c_i$ needs to be retained for the subsequent row insertions and for the final triangular solve $R_n w = c_n$. The scalar $\eta_i$ may be discarded unless the final residual norm is also required.

\noindent At step $k$, the current matrix $R$ is already upper triangular and the
previous components $a_1,\ldots,a_{k-1}$ of the inserted row have already
been annihilated. Therefore, in column $k$, the only nonzero entry below the
diagonal is the active element $a_k$.

\noindent For this reason, it is not necessary to build a Householder reflector on the
entire column. It is sufficient to consider only the two entries involved in
the elimination,

\begin{equation}
    z =
    \begin{bmatrix}
        R_{kk} \\
        a_k
    \end{bmatrix}
    \in \R^2.
\end{equation}

\noindent A two-dimensional normalized Householder reflector is then constructed as

\begin{equation}
    H_k^{(i)} = I_2 - 2u u^T,
    \qquad \norm{u}_2=1,
\end{equation}

\noindent and chosen so that

\begin{equation}
    H_k^{(i)}
    \begin{bmatrix}
        R_{kk} \\
        a_k
    \end{bmatrix}
    =
    \begin{bmatrix}
        \rho \\
        0
    \end{bmatrix}.
\end{equation}

\noindent Thus, at each step only the pair $(R_{kk},a_k)$ is needed to determine the
reflector: the diagonal entry is updated to $\rho$, while the corresponding
entry of the inserted row is annihilated.

\noindent The same $2\times2$ reflector must then be applied to the remaining entries
of the two affected rows, so that the transformation is consistent across
the whole trailing part of the matrix:

\begin{equation}
    \begin{bmatrix}
        R_{k,k:m} \\
        a_{k:m}
    \end{bmatrix}
    \leftarrow
    H_k^{(i)}
    \begin{bmatrix}
        R_{k,k:m} \\
        a_{k:m}
    \end{bmatrix}.
\end{equation}

\noindent The corresponding right-hand-side entries are transformed by the same reflector,

\begin{equation}
    \begin{bmatrix}
        c_k \\
        \eta
    \end{bmatrix}
    \leftarrow
    H_k^{(i)}
    \begin{bmatrix}
        c_k \\
        \eta
    \end{bmatrix},
\end{equation}

\noindent Here $\eta$ is initialized as $\eta=y_i$ at the beginning of the insertion of row $x_i^T$.

\noindent After repeating this procedure for $k=1,\ldots,m$, all entries of the inserted row have been annihilated. The row can then be discarded, while the updated $R$ and $c$ contain the effect of the newly inserted data row.

\begin{algorithm}[H]
\caption{Two-Dimensional Householder Row Insertion}
\label{alg:row_insertion_householder}
\begin{algorithmic}[1]
    \Require $X\in\R^{m\times n}$, $y\in\R^n$, $\lambda>0$
    \State $R \gets \lambda I_m$; \quad $c \gets 0\in\R^m$
    \For{$i=1,\ldots,n$}
        \State $a \gets X_{:,i}$; \quad $\eta \gets y_i$
        \For{$k=1,\ldots,m$}
            \State $z \gets [R_{kk},\, a_k]^{T}$
            \State $(u,\rho)\gets\Call{StableHouseholderVector}{z}$
            \State $H\gets I_2-2uu^T$
            \State $[R_{k,k:m};\,a_{k:m}] \gets H [R_{k,k:m};\,a_{k:m}]$
            \State $[c_k;\,\eta] \gets H[c_k;\,\eta]$
            \State Store the reflector record $(k,m+i,u)$
        \EndFor
    \EndFor
    \State Solve $Rw=c$ by back-substitution
    \State \Return $w$
\end{algorithmic}
\end{algorithm}

\subsection{Runtime Scalability}
\label{sec:qr_scalability}

\noindent The column-oriented structure-based method applies one reflector to an
$(n+1)\times(m-k+1)$ block at each step. Its factorization cost is therefore

\begin{equation}
\label{eq:structure_based_qr_cost}
    \sum_{k=1}^{m}O\bigl((n+1)(m-k+1)\bigr)
    =O(nm^2+m^2).
\end{equation}

\noindent Applying the reflectors to the right-hand side costs $O(nm)$ and back
substitution costs $O(m^2)$. The complete structure-based solver thus costs
$O(nm^2+m^2)$. The row-insertion algorithm has the same order:
for each of the $n$ data rows, it applies $m$ transformations to trailing
segments, at a cost

\begin{equation}
    \sum_{i=1}^{n}\sum_{k=1}^{m}O(m-k+1)=O(nm^2),
\end{equation}

\noindent followed by $O(m^2)$ back substitution.

\noindent Combining these bounds with Section~\ref{sec:qr_cost_project} gives

\begin{equation}
\label{eq:qr_scaling_summary}
\begin{aligned}
    T_{\mathrm{dense\ QR}}(m,n) &= O(m^3+nm^2),\\
    T_{\mathrm{structure\text{-}based\ QR}}(m,n) &= O(m^2+nm^2),\\
    T_{\mathrm{row\ insertion}}(m,n) &= O(m^2+nm^2),\\
    T_{\mathrm{normal\ equations}}(m,n) &= O(m^3+nm^2).
\end{aligned}
\end{equation}

\noindent Hence, with the project value $n=12$ fixed, both structured algorithms
are quadratic in $m$, whereas generic dense QR and a dense normal-equations
solve retain a cubic term. The structure-based implementation stores the
triangular factor, a copy of the dense block, and $m$ reflectors of length
$n+1$, for $O(m^2+nm)$ memory. The current row-insertion wrapper also stores
all $nm$ two-dimensional reflector records; disabling that optional storage
would leave only $O(m)$ workspace in addition to $R$.

\noindent The notebook first constructs a single matrix with $4000$ rows by
calling \texttt{build\_augmented\_X} with seed $42$ and jitter $0.01$. The
original $500$ rows are retained as the leading block; each additional row is
sampled as a complete row of $X$ and perturbed by Gaussian noise scaled by the
empirical standard deviation of each column. Taking prefixes of this matrix
produces nested problems with fixed $n=12$,
$m\in\{100,250,500,1000,2000,4000\}$, and $\lambda=0.5$. Data generation is
performed once, outside every timed region. Each entry in
Table~\ref{tab:qr_runtime_scaling} is the arithmetic mean of ten runs, without
a warm-up or randomized execution order.

\begin{table}[H]
\centering
\caption{Internal solver times in milliseconds for the augmented-data runtime
sweep. The normal-equations timer is partial.}
\label{tab:qr_runtime_scaling}
\small
\setlength{\tabcolsep}{5pt}
\begin{tabular}{rrrrr}
    \toprule
    $m$ & \textbf{Normal eq.} & \textbf{Dense QR}
        & \textbf{Structure-based QR} & \textbf{Row insertion}\\
    \midrule
     100 &    0.138 &      5.440 &   3.840 &  12.543\\
     250 &    1.333 &     70.165 &  18.151 &  68.379\\
     500 &    6.235 &    288.259 &  40.001 &  75.609\\
    1000 &   39.286 &   3830.648 &  61.597 & 143.880\\
    2000 &  296.157 &  26490.742 & 157.124 & 314.136\\
    4000 & 1628.571 & 180615.340 & 422.753 & 777.798\\
    \bottomrule
\end{tabular}
\end{table}

\noindent The solver timers still do not have identical scopes. The
normal-equations baseline measures only \texttt{numpy.linalg.solve}, excluding
construction of $XX^T+\lambda^2I_m$, construction of $Xy$, and evaluation of
the objective. In the current code, dense QR starts its timer before building
$A$ and $b$, while structure-based QR starts before building $[0;y]$; both
then include factorization, right-hand-side transformation, and back
substitution. Row insertion includes its fused right-hand-side updates and the
storage of all reflectors. Consequently, the normal-equations column prevents
Table~\ref{tab:qr_runtime_scaling} from being a uniform end-to-end comparison.

\noindent Within the comparable QR timings, at $m=4000$ the structure-based
path takes approximately $0.423$ seconds and row insertion $0.778$ seconds,
versus $180.615$ seconds for dense QR; the structure-based method is therefore
about $427$ times faster than dense QR at this size. On the final doubling from
$m=2000$ to $m=4000$, the dense, structure-based, and row-insertion times
increase by factors of approximately $6.82$, $2.69$, and $2.48$, respectively,
which is consistent with the cubic and quadratic terms in
Eq.~\eqref{eq:qr_scaling_summary}. These six measurements are empirical
evidence, not a proof of asymptotic complexity; moreover, the notebook computes
reference-slope arrays but does not draw them.

\subsection{Algorithmic Stability and the Condition Number}
\label{sec:qr_stability}

\noindent Conditioning is a property of the mathematical problem, whereas
stability is a property of the algorithm used to solve it. For
$A=[X^T;\lambda I_m]$,

\begin{equation}
    A^T A=XX^T+\lambda^2I_m.
\end{equation}

\noindent Since the project has $m>n$, the matrix $XX^T$ has at least $m-n$
zero eigenvalues. Therefore

\begin{equation}
\label{eq:qr_condition_lambda}
    \kappa_2(A)
    =\frac{\sqrt{\sigma_{\max}(X)^{2}+\lambda^2}}{\lambda},
    \qquad \lambda>0.
\end{equation}

\noindent If $X\neq0$, Eq.~\eqref{eq:qr_condition_lambda} behaves as
$\sigma_{\max}(X)/\lambda$ when $\lambda\to0^+$ and tends to $1$ as
$\lambda\to\infty$. Forming the normal equations squares this number, as
shown in Eq.~\eqref{eq:qr_condition_squaring}; this is precisely the
conditioning penalty avoided by Householder QR.

\noindent In the standard floating-point model with machine precision
$\varepsilon_{\mathrm{mach}}$, Householder QR is backward stable: up to
dimension-dependent constants, the computed factors are the exact QR factors
of $A+\Delta A$ with

\begin{equation}
    \frac{\norm{\Delta A}_2}{\norm{A}_2}=O(\varepsilon_{\mathrm{mach}}),
\end{equation}

\noindent and the computed orthogonal factor is orthogonal to working precision.
The resulting forward error still depends on the conditioning and geometry of
the least-squares problem. Accordingly, the notebook uses
$\varepsilon_{\mathrm{mach}}\kappa_2(A)$ for QR and
$\min\{1,\varepsilon_{\mathrm{mach}}\kappa_2(A)^{2}\}$ for normal equations
only as heuristic trend guides. They are not universal error bounds: the
forward condition number of the least-squares solution also depends on the
residual geometry. The notebook additionally plots
$\varepsilon_{\mathrm{mach}}$ as a reference scale for
$\eta_{\mathrm{opt}}$.

\noindent The notebook diagnoses this behavior using an independent thin-SVD
reference. If $X=U\Sigma V^T$, then the Tikhonov solution is

\begin{equation}
\label{eq:qr_svd_reference}
    w_{\mathrm{SVD}}
    =U\operatorname{diag}\!\left(
        \frac{\sigma_i}{\sigma_i^2+\lambda^2}
      \right)V^T y.
\end{equation}

\noindent For a computed solution $w$, the two reported diagnostics are the
relative forward error

\begin{equation}
\label{eq:qr_forward_error}
    e_{\mathrm{fwd}}
    =\frac{\norm{w-w_{\mathrm{SVD}}}_2}{\norm{w_{\mathrm{SVD}}}_2}
\end{equation}

\noindent and, for the augmented residual $r=Aw-b\neq0$, the scaled first-order
optimality residual

\begin{equation}
\label{eq:qr_scaled_optimality}
    \eta_{\mathrm{opt}}
    =\frac{\norm{A^T r}_2}{\norm{A}_2\norm{r}_2}.
\end{equation}

\noindent At an exact least-squares solution $A^T r=0$, so a small
$\eta_{\mathrm{opt}}$ verifies the optimality condition. It does not by itself
guarantee a small forward error when $A$ is severely ill-conditioned; both
diagnostics must therefore be considered together.

\noindent The experiment sweeps $\lambda\in\{10^j:j=-12,\ldots,4\}$. The
current dense solver receives $X$, $y$, and the current value of $\lambda$ at
every iteration, so its augmented system is now rebuilt correctly throughout
the sweep. For normal equations, dense QR, and structure-based QR,
respectively, the relative forward errors at $\lambda=10^{-6}$, where
$\kappa_2(A)\approx7.575\times10^7$, are approximately
$4.04\times10^{-1}$, $6.08\times10^{-14}$, and $1.11\times10^{-9}$. At
$\lambda=10^{-12}$, where
$\kappa_2(A)\approx7.575\times10^{13}$, they are approximately
$1.55\times10^1$, $4.04\times10^{-14}$, and $6.43\times10^{-4}$. At these two
ill-conditioned sample points, both Householder methods are substantially more
accurate than the explicitly formed normal equations, and dense QR has the
smallest forward error. Its near-machine-precision forward error at very small
$\lambda$ is an empirical property of this problem and must not be interpreted
as a universal conditioning-independent guarantee.

\noindent In the same method order, the corresponding
$\eta_{\mathrm{opt}}$ values are $6.63\times10^{-9}$,
$4.89\times10^{-10}$, and $9.03\times10^{-10}$ at $\lambda=10^{-6}$, and
$1.92\times10^{-3}$, $2.01\times10^{-3}$, and $6.64\times10^{-5}$ at
$\lambda=10^{-12}$. The extreme case also illustrates why forward error and
scaled optimality must be read together: dense QR has the smallest forward
error there, whereas structure-based QR has the smallest scaled optimality
residual.

\noindent One execution-state limitation remains in the saved notebook. Although
the row-insertion source has been corrected, the run retained the old imported
function in the active kernel. Its stored row-insertion accuracy and stability
values are therefore obsolete and are excluded from the conclusions above.
The printed scaled-optimality table correctly indexes the stored arrays, and
the condition numbers, SVD reference calculations, normal-equations results,
dense-QR results, and structure-based results are valid and support the
comparison made here.
